% Section 8. Requires amsmath, amssymb, amsthm, amsfonts.
% Uses the notation, theorem environments, and bibliography of Sections 2--6.
\section{Confidence Intervals for the VPC}
\label{sec:vpc_confidence}

Fix a covariate value $x$ and write $r_x(\theta)=\rho(x;\theta)$.
As in Section~\ref{sec:geometry_asymptotics}, let
\[
\theta=(\lambda^{\mathsf T},\gamma^{\mathsf T})^{\mathsf T}
\in\Theta=\Lambda\times K,\qquad
\gamma=(\gamma_{11},\gamma_{12},\gamma_{22})^{\mathsf T},
\]
where $\lambda$ collects all fixed-effect and family-specific parameters,
$K=\{\gamma:\Gamma(\gamma)\succeq0\}$, and
$d_\theta=p_\lambda+3$ is the dimension of the full parameter vector.
Let $n$ denote the number of independent clusters, and let $I_0\succ0$
be the limiting full Fisher information at the true parameter $\theta_0$.
The covariance matrix may be positive definite, nonzero rank one, or zero.

The local likelihood and constrained-estimation framework of
Section~\ref{sec:geometry_asymptotics} is maintained throughout.
In particular, its assumptions are high-level conditions on the marginal
likelihood, not restrictions to a particular conditional response family.
We distinguish fixed-parameter limits from assertions uniform over
sequences approaching changes in covariance rank, as in
Section~\ref{sec:pointwise_inference}.

For the strengthened compact-uniform statements below, let
$\mathcal K\subset\Theta$ be a compact true-parameter class buffered
from artificial nuisance and upper-covariance fitting boundaries,
while allowing all covariance ranks. The exact marginal likelihood
is globally maximized on a fixed enlarged compact fitted domain.
The four-family sufficient conditions require independent clusters,
a fixed finite collection of bounded informative designs with stable
positive frequencies, fixed finite nuisance dimension, and uniformly
positive-definite aggregate information. Gaussian residual variance,
NB2 size, and CPG dispersion remain bounded above and away from zero;
CPG power remains in a compact subinterval of $(1,2)$.
These conditions are verified in the completed likelihood supplement.
In particular, both CPG dispersion and power remain unknown.
The general fixed-parameter propositions below remain valid under
their individually stated assumptions, without asserting that these
stronger sufficient conditions are necessary.

\subsection{Interior VPC inference}
\label{subsec:ci-interior}

If $\Gamma_0\succ0$, $\lambda_0$ is interior, and $r_x$ admits a
continuously differentiable local extension, Theorem~\ref{thm:sec4-full-mle}
and the ordinary delta method yield
\begin{equation}\label{eq:sec8-interior}
\sqrt n\{r_x(\widehat\theta_n)-r_x(\theta_0)\}
\xrightarrow{d}N(0,\sigma_x^2),
\end{equation}
where
\begin{equation}\label{eq:sec8-interior-var}
\sigma_x^2=\nabla r_x(\theta_0)^{\mathsf T}
I_0^{-1}\nabla r_x(\theta_0).
\end{equation}
For $\sigma_x^2>0$, consistent estimation of the gradient and full
information gives the usual Wald interval, intersected with $[0,1]$.
The interior normal approximation is not automatically valid when
$\Gamma_0$ is singular.

\subsection{Functional inference under constrained estimation}
\label{subsec:ci-boundary}

Let $T_0=T_\Theta(\theta_0)$. Under the conditions of
Theorem~\ref{thm:sec4-full-mle},
\begin{equation}\label{eq:sec8-mle-limit}
\begin{aligned}
\sqrt n(\widehat\theta_n-\theta_0)&\xrightarrow{d}H_0,\\
H_0&=\Pi_{T_0}^{I_0}(I_0^{-1}Z),
\qquad Z\sim N_{d_\theta}(0,I_0).
\end{aligned}
\end{equation}
If $r_x$ is Hadamard directionally differentiable tangentially to $T_0$
under the conditions of Subsection~\ref{subsec:functional_asymptotics},
then
\begin{equation}\label{eq:sec8-directional}
\sqrt n\{r_x(\widehat\theta_n)-r_x(\theta_0)\}
\xrightarrow{d}V_x:=r'_{x,\theta_0}(H_0).
\end{equation}
For an ambient differentiable functional,
$r'_{x,\theta_0}(h)=\nabla r_x(\theta_0)^{\mathsf T}h$.
The limit can be non-Gaussian because the projection in
\eqref{eq:sec8-mle-limit} is generally nonlinear.
The projection acts on the full parameter vector. Its covariance component
may equivalently be obtained using the efficient information $J_0$ of
\eqref{eq:sec4-efficient-information}, followed by the nuisance adjustment
in \eqref{eq:sec4-nuisance-projection}. Thus the use of $I_0$ here and
$J_0$ in Sections~\ref{sec:global_lrt} and
\ref{sec:pointwise_inference} is consistent.

\subsection{Zero VPC at a nonzero rank-one covariance}
\label{subsec:ci-zero-rankone}

Set $z_x=(1,x)^{\mathsf T}$,
$d_x=(1,2x,x^2)^{\mathsf T}$, and
$v_x(\Gamma)=z_x^{\mathsf T}\Gamma z_x=d_x^{\mathsf T}\gamma$.
Suppose $\Gamma_0$ is nonzero rank one and $r_x(\theta_0)=0$.
Assume the observation-level variance is positive and, locally,
\begin{equation*}
\begin{aligned}
r_x(\theta)&=R_x(v_x(\Gamma),\lambda),\\
R_x(0,\lambda)&=0,\qquad
c_x=\partial_vR_x(0,\lambda_0)>0.
\end{aligned}
\end{equation*}
Put $g_x=(0_{p_\lambda}^{\mathsf T},d_x^{\mathsf T})^{\mathsf T}$.

\begin{proposition}[Positive-part Gaussian limit]
\label{prop:sec8-positive-part}
Under the preceding assumptions and Theorem~\ref{thm:sec4-full-mle},
\begin{equation}\label{eq:sec8-positive-part}
\sqrt n\,r_x(\widehat\theta_n)
\xrightarrow{d}s_x(N)_+,
\qquad N\sim N(0,1),
\end{equation}
where
\begin{equation}\label{eq:sec8-positive-scale}
s_x=c_x\sqrt{g_x^{\mathsf T}I_0^{-1}g_x}>0.
\end{equation}
The limiting distribution has an atom of probability $1/2$ at zero.
\end{proposition}

The proof is supplied in Section~7.1 of the completed confidence-set supplement. The full
inverse information incorporates nuisance-parameter correlations.


This reduction is specific to a zero-VPC functional at a rank-one
covariance. It is not a general assertion that smooth functionals at
rank-one points have positive-part Gaussian limits.

\subsection{Zero VPC at the covariance vertex}
\label{subsec:ci-zero-vertex}

At $\Gamma_0=0$, the tangent cone is
$T_0=\mathbb R^{p_\lambda}\times K$.
Under the same local representation and differentiability assumptions,
\begin{equation}\label{eq:sec8-vertex}
\sqrt n\,r_x(\widehat\theta_n)
\xrightarrow{d}
c_xg_x^{\mathsf T}
\Pi_{\mathbb R^{p_\lambda}\times K}^{I_0}(I_0^{-1}Z).
\end{equation}
The limiting variable is nonnegative but generally does not have the
positive-part Gaussian law in \eqref{eq:sec8-positive-part}.
Its distribution can be evaluated through Gaussian simulation and
Fisher-metric projection onto the PSD cone. In particular, the omnibus
chi-bar-square weights of Section~\ref{sec:global_lrt} do not directly
calibrate this scalar functional.

\subsection{Rank-adaptive projection inference}
\label{subsec:ci-rank-adaptive}

Consistency of $\widehat\Gamma_n$ does not imply convergence of its
unregularized tangent cone. At a vertex truth, for example, an estimated
positive-definite covariance has tangent cone $\mathbb R^3$, whereas
the true tangent cone is $K$. To estimate the fixed-parameter projection
law, let $\widehat a_1\ge\widehat a_2\ge0$ be the eigenvalues of
$\widehat\Gamma_n$, and choose a deterministic sequence $\tau_n$ with
\begin{equation}\label{eq:sec8-threshold}
\tau_n\downarrow0,\qquad\sqrt n\tau_n\longrightarrow\infty.
\end{equation}
Define $\widehat q_n=\sum_{j=1}^2\mathbf1\{\widehat a_j>\tau_n\}$.
Set $\widehat T_{K,n}=K$ when $\widehat q_n=0$,
$\widehat T_{K,n}=\mathbb R^3$ when $\widehat q_n=2$,
and otherwise use the supporting half-space determined by the
estimated kernel eigenvector. Put
$\widehat T_n=\mathbb R^{p_\lambda}\times\widehat T_{K,n}$.

\begin{proposition}[Fixed-parameter rank recovery]
\label{prop:sec8-rank-recovery}
If $\Gamma_0$ is fixed and
$\|\widehat\Gamma_n-\Gamma_0\|_{\mathrm{op}}=O_p(n^{-1/2})$,
then under \eqref{eq:sec8-threshold},
\begin{equation}\label{eq:sec8-rank-recovery}
\Pr\{\widehat q_n=\operatorname{rank}(\Gamma_0)\}\longrightarrow1.
\end{equation}
Moreover, the projections onto $\widehat T_{K,n}$ converge in
probability, uniformly for arguments in compact sets, to the
projections onto $T_K(\gamma_0)$ in any fixed positive-definite metric.
\end{proposition}

The proof is supplied in Section~7.2 of the completed confidence-set supplement.


Let $\widehat I_n\xrightarrow{p}I_0$ be positive definite.
Conditionally on the observations, draw
$Z^{(b)}\sim N_p(0,\widehat I_n)$ and compute
\begin{align}
H^{(b)}&=\Pi_{\widehat T_n}^{\widehat I_n}
(\widehat I_n^{-1}Z^{(b)}),\label{eq:sec8-sim-H}\\
V_x^{(b)}&=r'_{x,\widehat\theta_n}(H^{(b)}).
\label{eq:sec8-sim-V}
\end{align}
Under consistent derivative estimation, rank recovery, and accurate
projection computation, this estimates the fixed-parameter limiting law
in \eqref{eq:sec8-directional}. When the relevant limiting quantiles are uniquely identified, the
limiting distribution is continuous and strictly increasing at them,
and the estimated quantiles are consistent, an equal-tail interval is
\begin{equation}\label{eq:sec8-equal-tail}
\left[
\widehat r_x-\frac{\widehat q_{1-\alpha/2}}{\sqrt n},
\widehat r_x-\frac{\widehat q_{\alpha/2}}{\sqrt n}
\right]\cap[0,1],\qquad
\widehat r_x=r_x(\widehat\theta_n).
\end{equation}
Here $\widehat q_a$ denotes the estimated $a$-quantile of the
simulated limit, not the rank estimator $\widehat q_n$.
At atoms, including the atom in \eqref{eq:sec8-positive-part},
quantile continuity and exact equal-tail coverage cannot be presumed.
Acceptance regions must then specify boundary-mass conventions and
have separately verified coverage. For a discontinuous limit, one may
instead use estimated cutoffs
$a_n\le b_n$ for the centered estimation error, but coverage requires
control of the acceptance events, not merely weak convergence of the
simulated distribution. A sufficient condition is that the cutoffs
converge in probability to finite $a<b$ with
$\Pr(V_x=a)=\Pr(V_x=b)=0$ and
$\Pr(a\le V_x\le b)\ge1-\alpha$; then the inverted interval
$[\widehat r_x-b_n/\sqrt n,\widehat r_x-a_n/\sqrt n]
\cap[0,1]$ has limiting coverage at least $1-\alpha$.
If an endpoint lies at an atom, the inclusion probability must be
proved separately. Section~7.3 of the completed confidence-set supplement gives the precise
argument and illustrates the zero-VPC rank-one case.
The rank-recovery statement is pointwise: eigenvalues of order
$n^{-1/2}$ are not uniformly separable from zero.

\subsection{Rank-free confidence regions and score inversion}
\label{subsec:ci-rankfree}

Rank-adaptive simulation approximates the distribution of a
constrained estimator. A different strategy is to construct a
confidence region for the full parameter and project that region
through $r_x$. Both constructions below avoid covariance-rank
classification, although they rely on different regularity conditions.

\paragraph{A Fisher-metric ellipsoid centered at the constrained MLE.}
Let $\widehat I_n\xrightarrow{p}I_0\succ0$ and define
\begin{equation}\label{eq:sec8-ellipsoid}
\mathcal E_{n,1-\alpha}
=\left\{\theta\in\Theta:
 n(\theta-\widehat\theta_n)^{\mathsf T}
 \widehat I_n(\theta-\widehat\theta_n)
 \le\chi^2_{d_\theta,1-\alpha}\right\}.
\end{equation}
Its projected VPC set is $r_x(\mathcal E_{n,1-\alpha})$;
when extrema exist, the interval hull is
\begin{equation}\label{eq:sec8-ellipsoid-interval}
\mathcal I_{n,1-\alpha}^{\mathrm{ell}}(x)
=\left[
\inf_{\theta\in\mathcal E_{n,1-\alpha}}r_x(\theta),
\sup_{\theta\in\mathcal E_{n,1-\alpha}}r_x(\theta)
\right].
\end{equation}

\begin{proposition}[Fixed-parameter ellipsoid coverage]
\label{prop:sec8-ellipsoid-coverage}
Suppose the conditions of Theorem~\ref{thm:sec4-full-mle}
hold at a fixed $\theta_0$, and
$\widehat I_n\xrightarrow{p}I_0\succ0$. Then
\begin{equation}\label{eq:sec8-ellipsoid-coverage}
\liminf_{n\to\infty}
\Pr_{\theta_0}\{\theta_0\in\mathcal E_{n,1-\alpha}\}
\ge1-\alpha.
\end{equation}
Consequently, the projected set and its interval hull have
asymptotic coverage at least $1-\alpha$ for $r_x(\theta_0)$.
\end{proposition}

The proof is supplied in Section~7.4 of the completed confidence-set supplement.


The ellipsoid can be conservative because it uses the full-dimensional
$\chi_{d_\theta}^2$ cutoff. It is centered at the constrained MLE,
not at an unconstrained estimate that might violate the PSD constraint.

\begin{theorem}[Uniform rank-free ellipsoid coverage]
\label{thm:sec8-ellipsoid-uniform}
Suppose the uniform exact-likelihood conditions stated at the beginning
of this section hold on $\mathcal K$, and
$\widehat I_n-I(\theta)=o_{P_\theta}(1)$ uniformly over
$\theta\in\mathcal K$. Then
\begin{equation}\label{eq:sec8-ellipsoid-uniform}
\liminf_{n\to\infty}\inf_{\theta\in\mathcal K}
P_\theta\{\theta\in\mathcal E_{n,1-\alpha}\}\ge 1-\alpha.
\end{equation}
Its projected confidence sets and interval hulls have the same
uniform coverage lower bound, for every fixed $x$, including
sequences crossing covariance-rank boundaries.
\end{theorem}

The proof uses the \emph{true-centered} translated feasible set
$\mathbb R^{p_\lambda}\times\{\sqrt n(K-\gamma)\}$.
Its Gaussian quadratic maximizer is the metric projection of the
full score onto this closed convex set containing zero, and its
squared norm is bounded by the unrestricted full-score norm.
Uniform likelihood expansion, estimator localization, and the
$\chi^2_{d_\theta}$ score limit establish the result.
The full argument is given in Section~7.5 of the completed proof
supplement. No uniform covariance-rank estimator is used.

\paragraph{Full-parameter score inversion.}
The LAN assumption of Section~\ref{subsec:lan} is formulated only
for feasible local perturbations and does not itself assert existence
of a full ambient likelihood gradient at singular covariance matrices.
For the present construction, assume additionally that a suitable
ambient differentiable extension exists at candidate parameters.
Define
\begin{equation}\label{eq:sec8-score-stat}
\begin{aligned}
\Delta_n(\theta)&=n^{-1/2}\nabla_\theta\ell_n(\theta),\\
S_n(\theta)&=\Delta_n(\theta)^{\mathsf T}
\widehat I_n(\theta)^{-1}\Delta_n(\theta),
\end{aligned}
\end{equation}
where $\widehat I_n(\theta)$ estimates the \emph{full}
candidate-parameter information. For the verified score theorem below,
one may use the outer-product-of-scores estimator
\begin{equation}\label{eq:sec8-opg}
\widehat I_n^{\mathrm{OPG}}(\theta)
=\frac1n\sum_{i=1}^n s_i(\theta)s_i(\theta)^{\mathsf T},
\qquad s_i(\theta)=\nabla_\theta\ell_i(\theta).
\end{equation}
If the candidate information matrix is not positive definite,
set $S_n(\theta)=0$ as a conservative candidate-wise fallback;
this event has uniformly vanishing probability at a true parameter
on the verified informative classes.
For $u\in r_x(\Theta)$, put
\begin{equation}\label{eq:sec8-profile-score}
\begin{aligned}
\Theta_x(u)&=\{\theta\in\Theta:r_x(\theta)=u\},\\
S_{n,x}^{\mathrm{prof}}(u)&=
\inf_{\theta\in\Theta_x(u)}S_n(\theta).
\end{aligned}
\end{equation}
To avoid any ambiguity when the infimum is not attained, define
\begin{equation}\label{eq:sec8-score-acceptance}
\mathcal A_{n,1-\alpha}^{\mathrm{score}}
=\{\theta\in\Theta:S_n(\theta)\le\chi^2_{d_\theta,1-\alpha}\}
\end{equation}
and the exact projected confidence set
\begin{equation}\label{eq:sec8-score-ci}
\mathcal C_{n,1-\alpha}^{\mathrm{score}}(x)
=r_x(\mathcal A_{n,1-\alpha}^{\mathrm{score}}).
\end{equation}
If each relevant infimum is attained, this equals
$\{u:S_{n,x}^{\mathrm{prof}}(u)\le\chi^2_{d_\theta,1-\alpha}\}$.
Without attainment, the latter set may include additional boundary
values; the distinction does not affect the coverage lower bound below.

\begin{proposition}[Fixed-parameter score coverage]
\label{prop:sec8-rankfree}
Fix $\theta_0\in\Theta$. Suppose the score is defined at $\theta_0$
and
\[
\Delta_n(\theta_0)\xrightarrow{d}N_{d_\theta}(0,I_0),
\qquad
\widehat I_n(\theta_0)\xrightarrow{p}I_0\succ0.
\]
Then
\begin{equation}\label{eq:sec8-score-pointwise-coverage}
\liminf_{n\to\infty}\Pr_{\theta_0}
\{r_x(\theta_0)\in
\mathcal C_{n,1-\alpha}^{\mathrm{score}}(x)\}
\ge1-\alpha.
\end{equation}
No covariance-rank estimator is required.
\end{proposition}

The proof follows by full-parameter score convergence and deterministic projection; see Section~8.1 of the completed proof supplement.


\begin{proposition}[Conditional uniform score coverage]
\label{prop:sec8-rankfree-uniform}
Let $\Theta^\star\subseteq\Theta$ and suppose
\begin{equation}\label{eq:sec8-uniform-score-approx}
\begin{aligned}
\sup_{\theta\in\Theta^\star}\sup_{t\in\mathbb R}
\big|\Pr_\theta\{S_n(\theta)\le t\}\\
-F_{\chi_{d_\theta}^2}(t)\big|\longrightarrow0.
\end{aligned}
\end{equation}
Then
\begin{equation}\label{eq:sec8-uniform-score-coverage}
\liminf_{n\to\infty}\inf_{\theta\in\Theta^\star}
\Pr_\theta\{r_x(\theta)\in
\mathcal C_{n,1-\alpha}^{\mathrm{score}}(x)\}
\ge1-\alpha.
\end{equation}
\end{proposition}

The proof follows by the uniform approximation and set inclusion; see Section~8.1 of the completed proof supplement.

\begin{theorem}[Uniform score confidence region]
\label{thm:sec8-score-verified}
On $\mathcal K$ suppose the marginal cluster scores are centered,
possess uniformly bounded $(2+\delta)$ moments for some $\delta>0$,
and the aggregate information
$I_n(\theta)=n^{-1}\sum_i\operatorname{Var}_\theta\{s_i(\theta)\}$
has eigenvalues uniformly bounded above and away from zero.
Use \eqref{eq:sec8-opg} in \eqref{eq:sec8-score-stat}, with the
candidate-wise fallback above. Then
\begin{equation}\label{eq:sec8-score-verified}
\begin{aligned}
\sup_{\theta\in\mathcal K}\big|
P_\theta\{\theta\in\mathcal A_{n,1-\alpha}^{\mathrm{score}}\}\\
-(1-\alpha)\big|\longrightarrow0.
\end{aligned}
\end{equation}
Consequently, the exact projected score set and interval hull
cover $r_x(\theta)$ with probability at least $1-\alpha$
compact-uniformly. This conclusion applies to the four
specified GLMM families on their verified compact informative-design
classes, without covariance-rank classification.
\end{theorem}

The proof is a uniformly studentized triangular-array score CLT.
The candidate-parameter OPG estimator is consistent at the true
candidate by an entrywise law of large numbers using
$(2+\delta)$ score moments; no Ekvall--Bottai
individual-cluster information-comparison assumption is needed.
See Section~8.1 of the completed proof supplement and Appendix~L.
The score region may be inverted for a contrast as well as for a VPC.
In particular, projecting one full-parameter confidence region
simultaneously protects every well-defined VPC or difference
functional under the same coverage event.


When $\mathcal A_{n,1-\alpha}^{\mathrm{score}}$ is nonempty,
its VPC interval hull is
\begin{equation}\label{eq:sec8-score-interval}
\mathcal I_{n,1-\alpha}^{\mathrm{score}}(x)
=\left[
\inf_{\theta\in\mathcal A_{n,1-\alpha}^{\mathrm{score}}}r_x(\theta),
\sup_{\theta\in\mathcal A_{n,1-\alpha}^{\mathrm{score}}}r_x(\theta)
\right].
\end{equation}
It contains the exact projected set and inherits its coverage lower
bound. If the acceptance region is empty, no interval is produced by
this definition; a conservative fallback such as $[0,1]$ must be
specified if an interval is required for every sample. Both endpoints
involve constrained global optimization; local numerical optima alone
need not preserve the theoretical confidence region. A certified outer
approximation to the acceptance region is conservative, whereas an
inner approximation may under-cover.
Neither the score set nor the ellipsoid is guaranteed to be the shorter
of the two methods.


\subsection{Likelihood-free CPG moment confidence regions}
\label{subsec:ci-cpg-moments}

A separate rank-free construction is available for the CPG family
without evaluating its marginal likelihood, covariance score, or
Fisher information. Consider independent informative clusters, each
containing conditionally independent replicates $U_1,U_2,U_3$ at
$x=0$ and $V_1,V_2$ at $x=1$, with common CPG dispersion $\phi>0$
and unknown power $\kappa\in(1,2)$. Here the parameter is
$\theta=(\beta_0,\beta_1,\phi,\kappa,
\gamma_{11},\gamma_{12},\gamma_{22})$ of dimension seven.
The repeated observations at a given level have the same conditional
mean, family parameters, and random-effects design.
Define the observable moment vector
\begin{equation}\label{eq:sec8-cpg-moments}
\begin{aligned}
g(Y)=(&U_1,V_1,U_1U_2,V_1V_2,U_1V_1,\\
     &U_1^2-U_1U_2,\\
     &U_1^3-3U_1^2U_2+2U_1U_2U_3)^{\mathsf T},
\end{aligned}
\end{equation}
and write $m(\theta)=E_\theta g(Y)$ and
$\Omega(\theta)=\operatorname{Var}_\theta\{g(Y)\}$.
The seven identifying moments admit an explicit smooth inverse at
every covariance rank. In particular, using the notation of
Appendix~L, the contrast moments $A$ and $B_3$ satisfy
\begin{equation}\label{eq:sec8-cpg-inverse}
\frac{m_0B_3}{A^2}
=\kappa\exp\{(\kappa-1)^2v_0\},
\end{equation}
which is strictly increasing in $\kappa\in(1,2)$ for $v_0\ge0$.
Thus both unknown CPG family parameters are identified.

For $N$ independent clusters of this design, let
$\bar g_N=N^{-1}\sum_i g(Y_i)$ and let $\widehat\Omega_N$
be their empirical covariance matrix. When
$\widehat\Omega_N\succ0$, put
\begin{equation}\label{eq:sec8-cpg-stat}
Q_N(\theta)=N\{\bar g_N-m(\theta)\}^{\mathsf T}
\widehat\Omega_N^{-1}\{\bar g_N-m(\theta)\}
\end{equation}
and define
\begin{equation}\label{eq:sec8-cpg-region}
\mathcal A_{N,1-\alpha}^{\mathrm{mom}}
=\{\theta\in\mathcal K:Q_N(\theta)
\le\chi^2_{7,1-\alpha}\}.
\end{equation}
On the event $\widehat\Omega_N\not\succ0$, set this region equal
to the entire admissible compact class $\mathcal K$.

\begin{theorem}[Uniform CPG moment-region coverage]
\label{thm:sec8-cpg-moment}
For the preceding seven-parameter CPG experiment, assume the
compact admissible parameter class and exogenous informative
cluster design of Appendix~L. Then
\begin{equation}\label{eq:sec8-cpg-uniform}
\sup_{\theta\in\mathcal K}
\left|P_\theta\{\theta\in
\mathcal A_{N,1-\alpha}^{\mathrm{mom}}\}
-(1-\alpha)\right|\longrightarrow0.
\end{equation}
The projection of this region through $r_x$ or through any
fixed VPC contrast has uniform coverage at least $1-\alpha$.
\end{theorem}

The proof uses the smooth moment inverse, uniformly bounded
polynomial response moments, uniformly positive-definite
$\Omega(\theta)$, and a studentized multivariate central limit
theorem. These premises are verified in Appendix~L and in
Section~8.2 of the completed proof supplement.
This result is specific to the stated seven-parameter experiment;
it does not identify an arbitrary-dimensional nuisance vector from
only seven statistics. The verified grouped-design extension stacks
$D$ informative groups and uses $\chi^2_{7D,1-\alpha}$ with an
explicit full-rank design condition. A moment region may be less
efficient than a likelihood-based region, particularly because it
uses third-order response moments.

\subsection{Likelihood-ratio inversion and the pointwise test}
\label{subsec:ci-lr-inversion}

For $u\in r_x(\Theta)$, define the VPC-level likelihood ratio
\begin{equation}\label{eq:sec8-profile}
\mathrm{LR}_n(u)=2\left\{
\ell_n(\widehat\theta_n)-
\sup_{\theta\in\Theta_x(u)}\ell_n(\theta)
\right\}.
\end{equation}
Given critical values $c_{n,\alpha}(u)$, its inversion is
\begin{equation}\label{eq:sec8-inversion}
\mathcal C_{n,1-\alpha}^{\mathrm{LR}}(x)
=\{u\in r_x(\Theta):
\mathrm{LR}_n(u)\le c_{n,\alpha}(u)\}.
\end{equation}
The coverage of this set depends on valid size control for the
component tests, including nuisance parameters and boundary geometry.
At $u=0$, the level set is precisely the pointwise null ray of
Section~\ref{sec:pointwise_inference}, provided the zero-VPC
property holds. Thus its vertex and nonzero rank-one regimes, and the
local null sequences between them, are governed by the pointwise
analysis of that section rather than a universal $\chi_1^2$ cutoff.
The endpoint-envelope calibration developed there pertains to the
zero-VPC level set; it does not by itself calibrate the likelihood
ratio for arbitrary positive $u$. Nevertheless, a conservative
uniform all-level calibration is available without resolving the
geometry of every positive-level set.

\begin{theorem}[Uniform conservative LR inversion]
\label{thm:sec8-lr-universal}
Assume the compact-uniform exact-likelihood conditions on
$\mathcal K$ stated above: uniform root-$n$ localization of the full
constrained maximizer, true-centered feasible quadratic expansion,
uniform full-score Gaussian approximation, and uniformly
positive-definite information. Then for every exact null set
$A\subseteq\Theta$ containing the true $\theta$, the LRT
$\Lambda_{n,A}$, formed from the unrestricted and
null-restricted likelihood suprema, satisfies
\begin{equation}\label{eq:sec8-lr-universal}
\begin{aligned}
\limsup_{n\to\infty}\sup_{\theta\in\mathcal K}
\sup_{A\ni\theta}\ P_\theta\{\Lambda_{n,A}\\
>\chi^2_{d_\theta,1-\alpha}\}\le\alpha,
\end{aligned}
\end{equation}
where the inner supremum is over measurable exact-null
likelihood-ratio statistics. In particular, using the constant
cutoff $c_{n,\alpha}(u)=\chi^2_{d_\theta,1-\alpha}$ in
\eqref{eq:sec8-inversion} yields
\begin{equation}\label{eq:sec8-lr-uniform-coverage}
\liminf_{n\to\infty}\inf_{\theta\in\mathcal K}
P_\theta\{r_x(\theta)\in
\mathcal C_{n,1-\alpha}^{\mathrm{LR}}(x)\}
\ge1-\alpha.
\end{equation}
The same full-likelihood region provides conservative
simultaneous coverage for collections of VPCs or contrasts.
\end{theorem}

Indeed, for every $A\ni\theta$,
\begin{equation}\label{eq:sec8-lr-domination}
\begin{aligned}
0\le\Lambda_{n,A}
&\le2\{\ell_n(\widehat\theta_n)-\ell_n(\theta)\}\\
&\le\Delta_n(\theta)^{\mathsf T}
I(\theta)^{-1}\Delta_n(\theta)+o_{P_\theta}(1)
\end{aligned}
\end{equation}
uniformly on $\mathcal K$. The final quadratic form has the
full-dimensional $\chi^2_{d_\theta}$ limit. The crucial
bound is independent of the shape, convexity, or dimension
of $A$. Section~7.6 of the completed proof supplement gives
the complete uniform argument.

The cutoff is generally conservative, but it is valid at both
zero and positive VPC levels, across covariance-rank transitions.
The specialized endpoint cutoff of Section~\ref{sec:pointwise_inference}
can be substantially smaller at zero; it is not a valid default
at positive VPC levels. Smaller positive-level cutoffs require
additional level-set calibration and are not claimed here.

\subsection{Validity, computation, and scope}
\label{subsec:ci-scope}

The methods above answer different inferential questions.
The interior Wald interval is a regular fixed-parameter method;
the constrained functional delta method and rank-adaptive simulation
provide fixed-stratum limits, including non-Gaussian laws and atoms.
Rank recovery is not uniform over covariance-rank transitions.
The full-parameter ellipsoid, score-inversion region, and universal
profile-LR region provide conservative \emph{compact-uniform} coverage
under their respective verified likelihood and information premises.
The explicit CPG seven-moment construction offers a separate
likelihood-free uniform route on its informative compact model class.
None requires covariance-rank classification for uniform validity.

All projected VPC sets, difference sets, and simultaneous bands inherit
full-parameter coverage by deterministic inclusion. An interval hull
must contain the \emph{entire} image of the underlying confidence
region. Theoretical statements concern the exact confidence regions
and constrained global likelihood suprema. A local optimizer, fixed
quadrature tolerance, numerical score, or inner approximation to the
range is not enough to establish coverage. Certified outer bounds
preserve coverage, and the full-likelihood construction provides a
conservative alternative whenever a specialized critical value is
unavailable.

The compact-uniform four-family conclusions require the explicit
bounded informative-cluster designs, fixed nuisance dimension,
admissible buffered parameter compacts, uniform positive information,
and globally correct likelihood fitting described above. Neither
unbounded fitting domains, physical nuisance endpoints, singular
Fisher information, nor growing clusters are covered by those
verification theorems. The original mathematical results remain
valid under their individually stated premises outside the verified
classes, but further assumptions must then be established rather
than inferred from the family name alone. The completed confidence-set
proofs, including acceptance at atoms and both uniform projection
constructions, are in Sections~7--8 of the completed proof supplement;
Appendix~L contains the independent CPG moment-region derivation.
